\subsection{The bounded-type and one-sign checks behind C.28}
\begin{proposition}[Checking the projective-sequence classification C.21]
\label{F:C21-sequences}
For an actual projective sequence of types $t_n$, let $r_n$ count
its nonempty constant-sign runs. The following are equivalent:
$r_n$ is unbounded; every type receives a map from some $t_n$;
and the all-small-test formula in Proposition~\ref{F:C20-paths}
holds for every relative category and endpoint pair.
If these conditions fail, the tail of the path diagram is
isomorphic in $\Cat_{\mathrm{hot}}$ to the constant diagram at
one finite alternating type. The identifications commute with
the actual transition maps and localization effects.
\end{proposition}
\begin{proof}
In the allocation description of Lemma~\ref{F:type-contraction},
opposite successive target runs occupy distinct successive source
runs. Thus a type map cannot increase run count, and $r_n$ is
nondecreasing. Unbounded alternating run strings accommodate any
fixed target word, by allocating its successive letters to matching
runs and repeating positions within each run. Conversely, maps to
every alternating word require unbounded run counts. Theorem C.20,
checked in Proposition~\ref{F:C20-paths}, proves the formula in this
case with its actual comparison.

For necessity, if $r_n\le r$, take the finite alternating fence
with $r+2$ vertices, mark every arrow, and choose its two extremal
vertices as endpoints. A constant-sign path run moves by at most
one edge, since no two nonidentity arrows in either direction
compose. Thus all the path fibers are empty. In the groupoid
completion the two endpoints have exactly one arrow between them,
because the nerve is an interval. This is a complete finite
counterexample for every bounded run count, including $r=0$.

Finally bounded $r_n$ stabilizes. At a common positive run count,
a map of types must preserve the first sign: otherwise the target
run pattern needs one additional source run. Let $\sigma$ be this
fixed reduced alternating word. Selecting and repeating within
runs gives maps both ways between each tail type and $\sigma$.
Their composites and identities lie in the same nonempty mapping
posets, which contract by Lemma~\ref{F:type-contraction}.
Consequently their path pullbacks are inverse in
$\Cat_{\mathrm{hot}}$. The two maps to $\sigma$ associated to
each transition lie in that same mapping poset as well. This
checks compatibility with the given sequence, not only separate
equivalences of its stages. Applying
$\pi_0\operatorname{Fun}(X,-)$ preserves these finite uniform
zigzags for every small $X$. The zero-run case is the empty word.
The remaining endpoint-specific criterion is precisely that
every localization word has a $\sigma$ representative and equal
effects have connected $\sigma$ representatives; these are
surjectivity and injectivity of the specified effect map.
\end{proof}

\begin{proposition}[One-sign paths and the exact universal comparison]
\label{F:one-sign-paths}
For any positive length, the all-positive endpoint-fixed path
category is categorically homotopy equivalent to the discrete set
$\Hom_M(a,b)$. The all-negative one is equivalent to the discrete
set $\Sigma(b,a)$. These identifications commute with every
same-sign type pullback after passage to
$\Cat_{\mathrm{hot}}$, and hence hold after applying
$\pi_0\operatorname{Fun}(X,-)$ for every small $X$.

The comparison of positive paths with localization mapping spaces
is an equivalence at all endpoints exactly when every marked arrow
is already an isomorphism in $M$. For negative paths the analogous
all-endpoint condition is that $M$ is a groupoid and
$\Sigma=\operatorname{Mor}M$. At one endpoint pair the exact
condition is that the specified effect map is bijective on
components and every component of the localization mapping space
is contractible. In particular, when all arrows are marked, the
positive comparison holds everywhere exactly for groupoids.
\end{proposition}
\begin{proof}
Compose a positive path to its effect $f:a\to b$. A section sends
$f$ to the path with $f$ first and identities at $b$ afterwards.
The remaining composites from its intermediate objects to $b$
give a natural transformation from the original path to this
section. This is natural in endpoint-fixed ladders because each
square commutes. Thus composition and the section are inverse
up to a single transformation. For a negative path, use its
marked composite $b\to a$, put that arrow first and identities
after it, and use the marked composites from $b$ to each
intermediate object. They give the transformation in the reverse
direction. Closure of the marking ensures all these are paths.
Effect commutes strictly with type pullbacks. The sections then
give their claimed compatibility in $\Cat_{\mathrm{hot}}$.

The positive all-endpoint condition says exactly that the
localization functor $NM\to L(M,\Sigma)$ is fully faithful;
it is already surjective on objects. It is therefore an
equivalence, so it reflects equivalences and every member of
$\Sigma$ was invertible. Conversely localizing already invertible
arrows changes nothing, by the localization universal property.
For negative paths the structural functor
$\Sigma^{\mathrm{op}}\to L(M,\Sigma)$ sends each arrow to an
equivalence. If it is fully faithful, every arrow in $\Sigma$
has its inverse in $\Sigma$. These inverses are also inverses in
$M$, so the preceding argument identifies $L(M,\Sigma)$ with
$NM$. Full faithfulness now says every arrow of $M$ is the inverse
of one of $\Sigma$, proving the asserted necessity. The converse
and the single-endpoint criterion follow directly from the
discrete source mapping sets.

For the finite failure mentioned in C.6 one may use the four-arrow
circle category of Theorem~\ref{F:C17-fence}, with every arrow
marked. Its forward one-edge fiber at $(0,0)$ is a singleton,
whereas its localization is the groupoid completion of the
circle as a space, whose endpoint mapping space has components
$\mathbb Z$. The complete category table was checked there.
This replaces any reliance on the unavailable crown example.
\end{proof}

\paragraph{The composition assertions.}
For any fixed nonempty block $\tau$, paths in $\tau^n$ and
$\tau^m$ concatenate strictly to a path in $\tau^{n+m}$.
Padding by an identity block is a literal functor, so concatenation
defines the ind-operation with its units and associativity.
If the block has both signs, its run counts are unbounded and
Proposition~\ref{F:C21-sequences} identifies its all-small-test
component formula with localization. If it has one sign, the
preceding proposition identifies the operation with ordinary
composition (or inverse composition in the marked subcategory).
Choices of a common type give the same operation because the
mapping-poset contractions act uniformly on paths. These verify
the categorical and ind-operation statements used in C.28.
Their further identification with all higher localization
mapping-space compositions is exactly the remaining comparison
isolated in \ref{F:unmarked-path-comparison}; it is not inferred
from associativity on components.
